\fontsize{13pt}{22.5pt}\selectfont

\section{KẾT LUẬN}

\subsection{Kết quả đạt được}
\indent Đồ án đã hoàn thành các mục tiêu đề ra ban đầu. Kết quả chính là ba môi trường
huấn luyện game 3D có tính chất khác nhau trên nền tảng Unity ML-Agents: Football Table
(đối kháng một đối một, hành động liên tục, tự chơi), Capture The Flag (hợp tác hai tác
nhân với phần thưởng nhóm và chương trình học theo bậc) và Cross The Road (điều hướng
đơn tác nhân với phần thưởng thưa). Với mỗi môi trường, đề tài trình bày trọn vẹn quá
trình thiết kế không gian quan sát, không gian hành động và hàm phần thưởng, kể cả những
phương án đã thử và phải bỏ.

\indent Phần tri thức thiết kế thu được từ quá trình đó là kết quả đáng kể không kém
bản thân ba môi trường. Trường hợp điển hình là Capture The Flag: đặc tả ban đầu trừng
phạt đúng trạng thái phối hợp mà nhiệm vụ đòi hỏi - hai tác nhân đứng tách nhau ở hai
phía cửa - nên tác nhân mắc kẹt ở đáy phần thưởng nhóm bất kể thuật toán nào. Lỗi này
không biểu hiện thành lỗi chương trình mà chỉ thành một đường cong phẳng, và chỉ được
tìm ra khi quan sát hành vi thật. Bài học rút ra - hàm phần thưởng phải được kiểm tra
xem nó có vô tình phạt chính hành vi đích hay không - áp dụng được cho bất kỳ môi trường
hợp tác nào, không riêng môi trường này.

\indent Ba môi trường được đóng gói thành một gói Unity Package Manager kèm mười lăm
tệp cấu hình huấn luyện, bộ công cụ Editor và trình hướng dẫn cài đặt bảy bước, nên
chúng tồn tại được như một sản phẩm dùng lại chứ không chỉ như mã nguồn của một đồ án.
Đề tài cũng bổ sung hai thuật toán không có sẵn trong ML-Agents là MAPPO và DQN dưới
dạng trainer mở rộng, và hiện thực hóa quy trình kết nối môi trường Unity với một thư
viện học tăng cường bên ngoài qua Python Low-Level API.

\indent Mười lăm lần huấn luyện đã xác nhận cả ba môi trường đạt yêu cầu theo ba tiêu
chí đặt ra ở Chương~\ref{ch:results}. Chúng \textit{học được}: trên cả ba môi trường
đều có ít nhất bốn trong năm thuật toán tiến bộ rõ rệt so với lúc khởi tạo. Chúng
\textit{phân biệt được}: trên Cross The Road phần thưởng cuối trải từ $0{,}846$ (SAC)
xuống $0{,}336$ (DQN), còn trên Capture The Flag bốn thuật toán về đích trong dải
$0{,}800$--$0{,}840$ trong khi SAC sụp đổ ở $0{,}183$ - thứ hạng ở cả hai đều giải thích
được bằng cơ chế của bài toán. Và chúng \textit{không có lối tắt}: không lần chạy nào
đạt phần thưởng gần trần mà lại không giải được nhiệm vụ khi quan sát trực tiếp trong
ứng dụng trình diễn.

\indent Riêng Football Table cho thấy một giới hạn cần ghi nhận thẳng thắn: cả bốn thuật
toán ghép được với tự chơi đều tăng ELO từ ba mươi sáu tới sáu mươi bốn điểm so với
phiên bản cũ của chính mình, nghĩa là cơ chế tự chơi hoạt động đúng, nhưng biên độ đó
quá hẹp để môi trường phân định được các thuật toán với nhau. Ở cấu hình hiện tại,
Football Table đạt hai trong ba tiêu chí; muốn dùng nó làm thước đo so sánh thì cần
ngân sách bước lớn hơn hoặc một đối thủ mốc cố định thay cho tự chơi thuần túy.

\indent Một quan sát xuyên suốt cả ba môi trường ủng hộ chính luận điểm trung tâm của
đề tài: chênh lệch của cùng một thuật toán giữa hai môi trường luôn lớn hơn chênh lệch
giữa các thuật toán trong cùng một môi trường. Khoảng cách giữa thành công của SAC trên
Cross The Road và sụp đổ của chính nó trên Capture The Flag lớn hơn mọi khoảng cách nội
bộ. Nói cách khác, đặc tính bài toán - thứ do thiết kế môi trường quyết định - chi phối
kết quả mạnh hơn lựa chọn thuật toán, và đó là lý do khâu xây dựng môi trường xứng đáng
được đầu tư nhiều hơn mức thường thấy.

\indent Bên cạnh phần thực nghiệm so sánh, đề tài đã xây dựng hoàn chỉnh một ứng dụng
trình diễn tương tác, trong đó toàn bộ mười lăm mô hình đã huấn luyện được triển
khai trực tiếp vào game dưới định dạng ONNX. Ứng dụng cho phép người dùng lựa chọn môi
trường, lựa chọn một trong ba chế độ chơi (tự điều khiển, tác nhân tự chơi, chơi cùng
tác nhân), lựa chọn mô hình cụ thể cho tác nhân và
tra cứu ngay đường cong huấn luyện của cả năm thuật toán trên giao diện menu. Danh mục
mô hình cùng dữ liệu biểu đồ được sinh ra từ một kịch bản duy nhất chạy trên thư mục
kết quả, nên phần trình diễn và phần số liệu của Chương~\ref{ch:results} không thể lệch
nhau. Các luồng chức năng đã được kiểm thử trong Unity Editor và đều hoạt động đúng
(Mục~\ref{sub:app-testing}). Kết quả này chứng minh tính khả thi của
việc nhúng mô hình học tăng cường vào sản phẩm game hoàn chỉnh với chi phí suy luận
không đáng kể.

\subsection{Đóng góp của đề tài}
\indent Nếu phải xếp các đóng góp của đề tài theo mức độ hữu ích, tôi đặt bộ môi trường
lên đầu. Ba môi trường ba chiều bao phủ ba dạng bài toán học tăng cường phổ biến được
đóng gói thành một gói Unity Package Manager cài được vào dự án bất kỳ, kèm mười lăm cấu
hình huấn luyện đã chạy thật, bộ công cụ Editor và một trình hướng dẫn cài đặt bảy bước
tự kiểm tra trạng thái. Điều làm nên khác biệt so với việc chỉ công bố mã nguồn là ba
môi trường này đã được kiểm chứng: mỗi môi trường đã chịu năm thuật toán khác hẳn nhau
về nguyên lý, nên người dùng lại biết trước nó học được, nó phân biệt được, và nó không
có lối tắt - những điều mà một môi trường mới dựng không bao giờ bảo đảm được. Chính đồ
án này tiêu thụ gói đó qua ranh giới assembly đúng như một dự án bên ngoài, nên phần
trình diễn trở thành phép kiểm thử thường trực cho gói.

\indent Kế đến là phần tri thức thiết kế môi trường, thứ mà tài liệu về học tăng cường
thường bỏ qua. Đề tài ghi lại đầy đủ quá trình đặc tả cho ba dạng bài toán khác nhau,
bao gồm cả các phương án thất bại và dấu hiệu nhận biết chúng: một hàm phần thưởng
trừng phạt chính hành vi đích, một tác nhân thiếu quan sát về đồng đội nên không thể
định thời điểm phối hợp, một phạt thời gian chia cho không khi giới hạn bước bằng không.
Đi cùng là các quyết định vận hành có ảnh hưởng lớn nhưng ít được bàn tới - nhân khu vực
huấn luyện để tăng thông lượng, dựng biến thể không gian hành động rời rạc cho thuật
toán dựa trên giá trị - và hai con đường đưa một thuật toán không có sẵn vào Unity: gói
trainer mở rộng đăng ký MAPPO và DQN vào chính chương trình huấn luyện của ML-Agents, và
lớp bọc môi trường qua \texttt{mlagents\_envs} kết nối với Stable-Baselines3.

\indent Kết quả so sánh năm thuật toán, tuy là sản phẩm phụ của khâu kiểm chứng, vẫn có
giá trị tham khảo trực tiếp cho người làm game. Đáng chú ý nhất là một kết quả phủ định:
trên bài toán hợp tác có tác nhân đồng nhất, PPO thường bám sát cả MAPPO lẫn MA-POCA
trong dải chênh lệch không tới một phần trăm, nên chi phí tích hợp một thuật toán đa tác
nhân chuyên biệt chưa chắc được đền đáp. Đi cùng nó là một kết quả gây ngạc nhiên theo
chiều ngược lại: MAPPO vượt PPO một khoảng lớn trên chính môi trường đơn tác nhân, nơi
lý thuyết nói rằng phê bình tập trung của nó lẽ ra vô hiệu.

\indent Đóng góp còn lại là ứng dụng trình diễn với ba chế độ chơi, cơ chế chọn mô hình
cho phép đặt cạnh nhau cả mười lăm chính sách đã huấn luyện, và hai tab số liệu huấn
luyện nhúng thẳng trong game. Giá trị của nó không nằm
ở phần game, mà ở chỗ nó là công cụ soi lỗi đặc tả: quan sát hành vi thật là cách duy
nhất phát hiện những sai sót mà đường cong phần thưởng không thể hiện, và cũng là cách
khép kín quãng đường từ một tệp cấu hình YAML tới một hành vi
mà người chơi cảm nhận được bằng tay: cùng một bảng xếp hạng ở Chương~\ref{ch:results}
có thể được kiểm chứng lại bằng mắt, bằng cách cho hai mô hình khác thuật toán chạy
trên cùng một nhiệm vụ. Toàn bộ mã nguồn môi trường, cấu hình
huấn luyện và dữ liệu kết quả đều được công bố dưới dạng mở.

\subsection{Hạn chế}
\label{sub:limits}
\indent Hạn chế nặng nhất của đề tài là mỗi tổ hợp môi trường - thuật toán chỉ được chạy
với một seed duy nhất. Nghĩa là mọi chênh lệch báo cáo ở Chương~\ref{ch:results} đều
chưa tách được phần do bản chất thuật toán khỏi phần do may rủi của một lần khởi tạo.
Những khoảng cách lớn, như việc DQN xếp cuối trên Cross The Road hay việc SAC sụp đổ
trên Capture The Flag, nhiều khả năng vẫn đứng vững; nhưng các thứ hạng sát nhau thì
không nên đọc như
một kết luận thống kê. Đây là hệ quả của quỹ thời gian tính toán chứ không phải của một
lựa chọn phương pháp luận, và nếu chỉ được sửa một điều duy nhất trong đồ án này thì tôi
sẽ sửa đúng điều đó.

\indent Hạn chế kế tiếp nằm ở độ phân giải của phép đo. Ở Football Table, bốn thuật toán
tự chơi kết thúc trong dải ELO $1237$--$1260$; ở Capture The Flag, bốn thuật toán dẫn đầu
nằm trong dải phần thưởng $0{,}800$--$0{,}840$. Cả hai khoảng cách đó đều nhỏ hơn dao
động của chính các đường cong, nên hai môi trường này ở cấu hình hiện tại không phân biệt
được các thuật toán trong nhóm dẫn đầu. Chỉ Cross The Road cho một phổ đủ rộng để xếp
hạng có nghĩa. Muốn hai môi trường kia nói được nhiều hơn thì phải làm chúng khó lên, chứ
không phải chạy chúng lâu hơn.

\indent Một ô của lưới thực nghiệm cũng nằm ngoài phép so sánh: DQN trên Football Table
chạy không có tự chơi vì ML-Agents phát sinh lỗi khi ghép trainer off-policy với bộ điều
phối đối kháng. Lần chạy này vì thế không có ELO và học trong điều kiện khác hẳn bốn
thuật toán còn lại, nên mọi con số của nó ở môi trường này chỉ đọc được như một tham
chiếu rời.

\indent Sau cùng, đề tài không giải thích được hai hiện tượng mà chính số liệu của mình
phơi ra. Hiện tượng đầu là việc MAPPO vượt PPO trên môi trường đơn tác nhân, nơi hai
thuật toán chỉ khác nhau ở phần phê bình tập trung vốn lẽ ra vô hiệu khi nhóm chỉ có một
thành viên. Hiện tượng còn lại là kiểu thất bại của SAC trên Capture The Flag, khi
entropy của nó sụp xuống gần không rồi bật ngược trở lại đúng mức cực đại và ở đó tới
cuối. Cả hai đều cần thêm thực nghiệm mới kết luận được, và với một hạt giống cho mỗi ô
thì cũng chưa loại trừ được khả năng hiện tượng thứ nhất chỉ là dao động ngẫu nhiên.

\subsection{Hướng phát triển}
\indent Việc cần làm trước tiên là chạy lại toàn bộ mười lăm tổ hợp với nhiều hạt giống,
vì chừng nào chưa có phương sai giữa các lần chạy thì mọi so sánh trong đồ án vẫn còn
treo lơ lửng. Song song với đó là làm hai môi trường Football Table và Capture The Flag
khó lên, bởi ở cấu hình hiện tại chúng không đủ độ phân giải để tách các thuật toán dẫn
đầu; với Capture The Flag, hướng cụ thể là thêm bài trong chương trình học hoặc buộc hai
tác nhân phối hợp chặt hơn thay vì kéo dài thời lượng huấn luyện.

\indent Ngay sau đó là hai câu hỏi mà đề tài đặt ra nhưng chưa trả lời được. Câu hỏi thứ
nhất là vì sao MAPPO vượt PPO trên môi trường đơn tác nhân; một thực nghiệm tách riêng
phần phê bình tập trung, chạy trên nhiều hạt giống, sẽ cho biết đây là một hiệu ứng thật
hay chỉ là may rủi của một lần khởi tạo. Câu hỏi còn lại là cơ chế đứng sau kiểu thất bại
của SAC trên bài toán hợp tác; giả thuyết phù hợp với dữ liệu là bộ đệm phát lại lưu lẫn
lộn kinh nghiệm của hai bài học trong chương trình học, và cách kiểm chứng trực tiếp là
chạy lại với bộ đệm tách riêng theo từng bài. Ranh giới giữa chia sẻ tham số và attention
cũng đáng khảo sát bằng một môi trường hợp tác có hai tác nhân vai trò bất đối xứng, hoặc
có tác nhân rời nhóm giữa lượt chơi.

\indent Ở tầm xa hơn, bộ so sánh có thể được mở rộng bằng các thuật toán như DreamerV3
hay QMIX - việc bổ sung nay đã rẻ hơn nhiều nhờ cơ chế trainer mở rộng đã dựng sẵn - còn
ba môi trường có thể được nâng độ phức tạp lên gần với game thương mại.
Ứng dụng trình diễn thì cần một bản chạy độc lập (WebGL hoặc desktop), kèm bảng xếp
hạng người chơi đấu với tác nhân và một khảo sát trải nghiệm người dùng - đây là cách
duy nhất để trả lời câu hỏi mà không chỉ số huấn luyện nào chạm tới được, rằng hành vi
của tác nhân có thực sự ``tự nhiên'' trong mắt người chơi hay không.
